\documentclass{ximera}

\begin{document}
	\author{Stitz-Zeager}
	\xmtitle{Exercises for The Other Circular Functions}{}

\mfpicnumber{1} \opengraphsfile{ExercisesforTheOtherCircularFunctions} % mfpic settings added 

\begin{question}
In Exercises \ref{circvaluefirst} - \ref{circvaluelast}, find the exact value or state that it is undefined.

\begin{problem} \label{circvaluefirst}
$\tan \left( \frac{\pi}{4} \right)$
\end{problem}

\begin{problem}
$\sec \left( \frac{\pi}{6} \right)$
\end{problem}

\begin{problem}
$\csc \left( \frac{5\pi}{6} \right)$
\end{problem}

\begin{problem}
$\cot \left( \frac{4\pi}{3} \right)$
\end{problem}

\begin{problem}
$\tan \left( -\frac{11\pi}{6} \right)$
\end{problem}

\begin{problem}
 $\sec \left( -\frac{3\pi}{2} \right)$
\end{problem}
 
\begin{problem}
$\csc \left( -\frac{\pi}{3} \right)$
\end{problem}

\begin{problem}
$\cot \left( \frac{13\pi}{2} \right)$
\end{problem}

\begin{problem}
$\tan \left( 117\pi \right)$
\end{problem}

\begin{problem}
$\sec \left( -\frac{5\pi}{3} \right)$
\end{problem}

\begin{problem}
$\csc \left( 3\pi \right)$
\end{problem}


\begin{problem}
$\cot \left( -5\pi \right)$
\end{problem}

\begin{problem}
$\tan \left( \frac{31\pi}{2} \right)$
\end{problem}

\begin{problem}
$\sec \left( \frac{\pi}{4} \right)$
\end{problem}

\begin{problem}
$\csc \left( -\frac{7\pi}{4} \right)$
\end{problem}

\begin{problem}
$\cot \left( \frac{7\pi}{6} \right)$
\end{problem}

\begin{problem}
$\tan \left( \frac{2\pi}{3} \right)$
\end{problem}

\begin{problem}
$\sec \left( -7\pi \right)$
\end{problem}

\begin{problem}
$\csc \left( \frac{\pi}{2} \right)$
\end{problem}

\begin{problem} \label{circvaluelast}
$\cot \left( \frac{3\pi}{4} \right)$
\end{problem}

\end{question}

\begin{question}
In Exercises \ref{whereisanglefirst} - \ref{whereisanglelast}, use the given the information to determine the quadrant in which the terminal side of the angle lies when plotted in standard position.

\begin{problem} \label{whereisanglefirst} 
$\sin(\theta) > 0$ but $\tan(\theta) < 0$.
\end{problem}

\begin{problem}
$\cot(\alpha) > 0$ but $\cos(\alpha) < 0$.
\end{problem}

\begin{problem}
$\sin(\beta) > 0$ and $\tan(\beta) > 0$.
\end{problem}

\begin{problem} \label{whereisanglelast} $\cos(\gamma) > 0$ but $\cot(\gamma) < 0$.
\end{problem}

\end{question}

\begin{question}
In Exercises \ref{findothercircfirst} - \ref{findothercirclast}, use the given the information to find the exact values of the circular functions of $\theta$.

\begin{problem} \label{findothercircfirst}
$\sin(\theta) = \frac{3}{5}$ with $\theta$ in Quadrant II
\end{problem}

\begin{problem} 
$\tan(\theta) = \frac{12}{5}$ with $\theta$ in Quadrant III
\end{problem}

\begin{problem}
$\csc(\theta) = \frac{25}{24}$ with $\theta$ in Quadrant I
\end{problem}

\begin{problem}
$\sec(\theta) = 7$ with $\theta$ in Quadrant IV \end{problem}

\begin{problem}
$\csc(\theta) = -\frac{10\sqrt{91}}{91}$ with $\theta$ in Quadrant III
\end{problem}

\begin{problem}
$\cot(\theta) = -23$ with $\theta$ in Quadrant II
\end{problem}

\begin{problem}
$\tan(\theta) = -2$ with $\theta$ in Quadrant IV.
\end{problem}

\begin{problem}
$\sec(\theta) = -4$ with $\theta$ in Quadrant II.
\end{problem}

\begin{problem}
$\cot(\theta) = \sqrt{5}$ with $\theta$ in Quadrant III. \end{problem}

\begin{problem}
$\cos(\theta) = \frac{1}{3}$ with $\theta$ in Quadrant I.
\end{problem}

\begin{problem}
$\cot(\theta) = 2$ with $0  < \theta < \frac{\pi}{2}$.
\end{problem}

\begin{problem}
$\csc(\theta) = 5$ with $\frac{\pi}{2} < \theta < \pi$.
\end{problem}

\begin{problem}
$\tan(\theta) = \sqrt{10}$ with $\pi < \theta < \frac{3\pi}{2}$.
\end{problem}

\begin{problem} \label{findothercirclast}
$\sec(\theta) = 2\sqrt{5}$ with $\frac{3\pi}{2} < \theta < 2\pi$.
\end{problem}

\end{question}

\begin{question}
In Exercises \ref{circcalcfirst} - \ref{circcalclast}, use your calculator to approximate the given value to three decimal places.  Make sure your calculator is in the proper angle measurement mode!

\begin{problem} \label{circcalcfirst}
$\csc(78.95^{\circ})$
\end{problem}

\begin{problem}
$\tan(-2.01)$
\end{problem}

\begin{problem}
$\cot(392.994)$
\end{problem}

\begin{problem}
$\sec(207^{\circ})$
\end{problem}

 \begin{problem}
 $\csc(5.902)$
 \end{problem}
 
\begin{problem}
$\tan(39.672^{\circ})$
\end{problem}

\begin{problem}
$\cot(3^{\circ})$
\end{problem}

\begin{problem} \label{circcalclast}
$\sec(0.45)$ 
\end{problem}

\end{question}

\begin{question}

In Exercises \ref{circequanglefirst} - \ref{circequanglelast}, find all of the angles which satisfy the equation.

\begin{problem} \label{circequanglefirst}
$\tan(\theta) = \sqrt{3}$ 
\end{problem}

\begin{problem}
$\sec(\theta) = 2$
\end{problem}

\begin{problem}
$\csc(\theta) = -1$
\end{problem}

\begin{problem}
$\cot(\theta) = \frac{\sqrt{3}}{3}$
\end{problem}

\begin{problem}
$\tan(\theta) = 0$
\end{problem}

\begin{problem}
$\sec(\theta) = 1$
\end{problem}

\begin{problem}
$\csc(\theta) = 2$
\end{problem}

\begin{problem}
$\cot(\theta) = 0$
\end{problem}

\begin{problem}
$\tan(\theta) = -1$
\end{problem}

\begin{problem}
$\sec(\theta) = 0$
\end{problem}

\begin{problem}
$\csc(\theta) = -\frac{1}{2}$
\end{problem}

\begin{problem}
$\sec(\theta) = -1$
\end{problem}

\begin{problem}
$\tan(\theta) = -\sqrt{3}$
\end{problem}

\begin{problem}
$\csc(\theta) = -2$
\end{problem}

\begin{problem} \label{circequanglelast}
$\cot(\theta) = -1$ 
\end{problem}

\end{question}

\begin{question}
In Exercises \ref{circequtfirst} - \ref{circequtlast}, solve the equation for $t$.  Give exact values.

\begin{problem} \label{circequtfirst}
$\cot(t) = 1$
\end{problem}

\begin{problem}
$\tan(t) = \frac{\sqrt{3}}{3}$
\end{problem}

\begin{problem}
$\sec(t) = -\frac{2\sqrt{3}}{3}$
\end{problem}

\begin{problem}
$\csc(t) = 0$
\end{problem}

\begin{problem}
$\cot(t) = -\sqrt{3}$
\end{problem}

\begin{problem}
$\tan(t) = -\frac{\sqrt{3}}{3}$
\end{problem}

\begin{problem}
$\sec(t) = \frac{2\sqrt{3}}{3}$
\end{problem}

\begin{problem} \label{circequtlast}
$\csc(t) = \frac{2\sqrt{3}}{3}$ 
\end{problem}

\end{question}

\begin{question}
In Exercises \ref{decomposebasicothercircularfirst} - \ref{decomposebasicothercircularlast}, write the given function as a nontrivial decomposition of functions as directed.

\begin{problem} \label{decomposebasicothercircularfirst}
For $f(t) = 3t^2 + 2 \tan(3t)$, find functions $g$ and $h$ so that $f=g+h$. 
\end{problem}

\begin{problem}
For $f(\theta) = \sec(\theta) - \tan(\theta)$, find functions $g$ and $h$ so that $f=g-h$.
\end{problem}

\begin{problem}
For $f(t) = -\csc(t) \cot(t)$, find functions $g$ and $h$ so that $f=gh$.
\end{problem}

\begin{problem}
For $r(t) = \frac{\tan(3t)}{t}$, find functions $f$ and $g$ so $r = \frac{f}{g}$.
\end{problem}

\begin{problem}
For $T(\theta) =\tan(4 \theta)$, find functions $f$ and $g$ so $T = g \circ f$.
\end{problem}

\begin{problem}
For $s(\theta) = \sec^{2}(\theta)$, find functions $f$ and $g$ so $s = g \circ f$.
\end{problem}

\begin{problem}
For $L(x) = \ln (\sin(x) )$, find functions $f$ and $g$ so $L = g \circ f$.
\end{problem}

\begin{problem} \label{decomposebasicothercircularlast}
For $\ell(\theta) = \ln | \sec(\theta) - \tan(\theta)|$, find  find functions $f$,  $g$, and $h$ so $\ell = h \circ (f-g)$.
\end{problem}

\end{question}


\begin{problem} 
Let $S(t) = \sin(t)$ and $C(t) = \cos(t)$, $F(t) = \tan(t)$, and $G(t) = \cot(t)$.  Explain why $F = \frac{S}{C}$ but $F \neq \frac{1}{G}$.

HINT: Think about domains \ldots

\end{problem}

\begin{problem} \label{tangentarcexercise}
For each function $T(t)$ listed below, compute the average rate of change over the indicated interval.\footnote{See Definition \ref{arc} in Section \ref{AverageRateofChange} for a review of this concept, as needed.}  What trends do you notice? Compare your answer with what you discovered in Section \ref{TheCircularFunctionsSineandCosine} number \ref{sinearcexercise}.  Be sure your calculator is in radian mode!

\[ \begin{array}{|r||c|c|c|}  \hline

 T(t) &  [-0.1, 0.1] & [-0.01, 0.01] &[-0.001, 0.001] \\ \hline
 \tan(t)     &&&  \\  \hline
 \tan(2t)   &&&  \\ \hline
 \tan(3t)   &&&   \\  \hline
 \tan(4t)   &&&   \\  \hline

\end{array} \]

\end{problem}



\begin{problem} \label{sintovertexercise1} 
We wish to establish the inequality $\cos(\theta) < \frac{\sin(\theta)}{\theta} < 1$ for $0 < \theta < \frac{\pi}{2}.$  Use the diagram from the beginning of the section, partially reproduced below, to answer the following.

\begin{center}

\begin{mfpic}[20]{-1}{4}{-1}{6}
\axes
\drawcolor[gray]{0.7}
\parafcn{0,90,5}{3*dir(t)}
\drawcolor{black}
\arrow \parafcn{5, 55, 5}{0.75*dir(t)}
\tlabel[cc](0.75,0.5){\scriptsize $\theta$}
\point[3pt]{(0,0), (3,5.196), (3,0)}
\tlabel(3.75,-0.5){\scriptsize $x$}
\tlabel(0.25,5.75){\scriptsize $y$}
\tlabel(0.25,3.1){\scriptsize $1$}
\tlabel(-0.5,-0.5){\scriptsize $O$}
\tlabel(2,-0.5){\scriptsize $B(1,0)$}
\xmarks{0 step 3 until 3}
\ymarks{0 step 3 until 3}
\polyline{(0,0), (3,5.196)}
\polyline{(3,0), (3,5.196)}
\polyline{(2.75,0), (2.75, 0.25), (3,0.25)}
\polyline{(3,0), (1.5, 2.5981)}
\tlabel(1.75,2.6){\scriptsize $P$}
\tlabel(3.25,5.25){\scriptsize $Q$}
\end{mfpic} 

\end{center}

\begin{enumerate}

\item Show that triangle $OPB$ has area $\frac{1}{2} \sin(\theta)$ and triangle $OQB$ has area $\frac{1}{2} \tan(\theta)$.
\item Show that the circular sector $OPB$ with central angle $\theta$ has area $\frac{1}{2} \theta$.
\item Comparing areas, show that $\sin(\theta) < \theta < \tan(\theta)$ for $0 < \theta < \frac{\pi}{2}.$ 
\item Use the inequality $\sin(\theta) < \theta$ to show that $\frac{\sin(\theta)}{\theta} < 1$ for $0 < \theta < \frac{\pi}{2}.$
\item Use the inequality $\theta < \tan(\theta)$ to show that $\cos(\theta) < \frac{\sin(\theta)}{\theta}$ for $0 < \theta < \frac{\pi}{2}.$  Combine this with the previous part to complete the proof.

\end{enumerate}
\end{problem}

\begin{problem} \label{sintovertexercise2} 
Show that $\cos(\theta) < \frac{\sin(\theta)}{\theta} < 1$ also holds for $-\frac{\pi}{2}< \theta < 0$.
\end{problem}

\begin{problem}\label{sintovertexercise3}  
Use the results from Exercises \ref{sintovertexercise1}  and \ref{sintovertexercise2} along with the Squeeze Theorem,\footnote{Theorem \ref{squeezeth}}  to prove  $\ds{ \lim_{\theta \rightarrow 0}}$ $ \frac{\sin(\theta)}{\theta} = 1$.

\end{problem}

\end{document}

